%% ma: L11-B04-0006
%% lop: 11
%% chuong: 1
%% bai: 4
%% dang: 11.4.6
%% loai: TN4
%% muc: TH
%% dap_an: "C"
%% nguon: "de-cuong-giua-ky-1-toan-11-nam-2026-2027-truong-thpt-hoang-van-thu-ha-noi-2-8.pdf (M08, Đề cương giữa kì 1 Toán 11 THPT Hoàng Văn Thụ, Hà Nội), Phần 1 Câu 15"
%% kien_thuc: "Bài 4 (phương trình lượng giác cơ bản)"
%% trang_thai: chinh-thuc
%% ghi_chu: "Đã duyệt 29/10/2026."
\begin{ex}
Tập nghiệm $S$ của phương trình $\sqrt3\tan\dfrac{x}{3}+3=0$ là
\choice{$S=\left\{-\dfrac{\pi}{9}+k3\pi,\ k\in\mathbb Z\right\}$}{$S=\left\{-\dfrac{\pi}{3}+k\pi,\ k\in\mathbb Z\right\}$}{\True $S=\{-\pi+k3\pi,\ k\in\mathbb Z\}$}{$S=\left\{\dfrac{\pi}{6}+k\pi,\ k\in\mathbb Z\right\}$}
\loigiai{$\tan\dfrac x3=-\sqrt3\Leftrightarrow\dfrac x3=-\dfrac\pi3+k\pi\Leftrightarrow x=-\pi+k3\pi$ $(k\in\mathbb Z)$.}
\end{ex}
